\documentclass[11pt,a4paper]{article}

% ============================================================
%  MEDCONTROL — Documentação Técnica de Produto
%  Farmácia Inteligente | Relatório de Projeto
% ============================================================

\usepackage[a4paper,margin=2.2cm,top=2.8cm,bottom=2.8cm]{geometry}
\usepackage[utf8]{inputenc}
\usepackage[T1]{fontenc}
\usepackage[portuguese]{babel}
\usepackage{xcolor}
\usepackage{amssymb}
\usepackage{graphicx}
\usepackage{tikz}
\usetikzlibrary{shapes.geometric,arrows.meta,positioning,shadows,calc}
\usepackage{fancyhdr}
\usepackage{titlesec}
\usepackage{tcolorbox}
\tcbuselibrary{skins,breakable}
\usepackage{enumitem}
\usepackage{booktabs}
\usepackage{array}
\usepackage{longtable}
\usepackage{hyperref}
\usepackage{listings}
\usepackage{float}
\usepackage[section]{placeins}

% Ícones leves (substituem fontawesome5, evitando dependências indisponíveis)
\newcommand{\iInfo}{\textbf{\textcolor{white}{(i)}}}
\newcommand{\iWarn}{\textbf{\textcolor{white}{!}}}
\newcommand{\iChip}{\textbf{\textcolor{mcTealDark}{\#}}}
\newcommand{\iCheck}{\textbf{\textcolor{white}{\checkmark}}}
\newcommand{\iImage}{\textcolor{mcGray}{$\square$}}
\newcommand{\iGit}{\textbf{\textcolor{mcTealDark}{\textless/\textgreater}}}
\usepackage{multicol}
\usepackage{parskip}
\usepackage{tabularx}
\usepackage{makecell}
\usepackage{multirow}

% ------------------------------------------------------------
% Paleta de cores do produto — tema "saúde / farmácia inteligente"
% ------------------------------------------------------------
\definecolor{mcTeal}{HTML}{0E7C7B}      % cor primária
\definecolor{mcTealDark}{HTML}{0A5C5B}  % primária escura
\definecolor{mcMint}{HTML}{17BEBB}      % destaque
\definecolor{mcCoral}{HTML}{FF6B6B}     % alerta / atenção
\definecolor{mcAmber}{HTML}{F2A65A}     % aviso
\definecolor{mcGray}{HTML}{4A5568}      % texto secundário
\definecolor{mcLight}{HTML}{F4FAF9}     % fundo claro
\definecolor{mcLine}{HTML}{D8E8E6}      % linhas/bordas suaves

% ------------------------------------------------------------
% Hyperlinks
% ------------------------------------------------------------
\hypersetup{
    colorlinks=true,
    linkcolor=mcTealDark,
    urlcolor=mcTeal,
    citecolor=mcTeal,
    pdftitle={MedControl - Documentacao do Projeto},
    pdfauthor={Equipe MedControl},
    pdfsubject={Farmacia Inteligente - Relatorio de Projeto}
}

% ------------------------------------------------------------
% Cabeçalho / Rodapé estilo "documentação de produto"
% ------------------------------------------------------------
\pagestyle{fancy}
\setlength{\headheight}{23pt}
\fancyhf{}
\renewcommand{\headrulewidth}{0.8pt}
\renewcommand{\footrulewidth}{0.4pt}
\fancyheadoffset{0pt}
\renewcommand{\headrule}{\hbox to\headwidth{\color{mcLine}\leaders\hrule height \headrulewidth\hfill}}
\renewcommand{\footrule}{\hbox to\headwidth{\color{mcLine}\leaders\hrule height \footrulewidth\hfill}}

\fancyhead[L]{\small\color{mcGray}\textbf{MedControl}\,\textcolor{mcLine}{|}\,Farmácia Inteligente}
\fancyhead[R]{\small\color{mcGray}\leftmark}
\fancyfoot[L]{\small\color{mcGray}Documentação Técnica do Projeto}
\fancyfoot[R]{\small\color{mcGray}Página \thepage\ de \pageref{LastPage}}
\usepackage{lastpage}

% ------------------------------------------------------------
% Estilo de seções — visual "documentação"
% ------------------------------------------------------------
\titleformat{\section}
  {\Large\bfseries\color{mcTealDark}}
  {\colorbox{mcTeal}{\makebox[1.1cm]{\color{white}\thesection}}}
  {0.4em}
  {}
\titlespacing*{\section}{0pt}{1.6em}{1em}

\titleformat{\subsection}
  {\large\bfseries\color{mcTealDark}}
  {\thesubsection}
  {0.6em}
  {}
\titlespacing*{\subsection}{0pt}{1.2em}{0.6em}

\titleformat{\subsubsection}
  {\normalsize\bfseries\color{mcGray}}
  {\thesubsubsection}
  {0.5em}
  {}

% ------------------------------------------------------------
% Caixas estilo "documentação de produto" (nota, dica, especificação)
% ------------------------------------------------------------
\newtcolorbox{infobox}[1][]{
  enhanced, breakable, colback=mcLight, colframe=mcTeal,
  boxrule=0.6pt, left=8pt, right=8pt, top=6pt, bottom=6pt,
  arc=2pt, boxsep=2pt,
  title={\iInfo\ #1}, fonttitle=\bfseries\color{white},
  coltitle=white, colbacktitle=mcTeal,
  attach boxed title to top left={yshift=-2mm, xshift=4mm},
  boxed title style={boxrule=0pt, arc=2pt, colback=mcTeal}
}

\newtcolorbox{warnbox}[1][]{
  enhanced, breakable, colback=mcLight, colframe=mcAmber,
  boxrule=0.6pt, left=8pt, right=8pt, top=6pt, bottom=6pt,
  arc=2pt, boxsep=2pt,
  title={\iWarn\ #1}, fonttitle=\bfseries\color{white},
  coltitle=white, colbacktitle=mcAmber,
  attach boxed title to top left={yshift=-2mm, xshift=4mm},
  boxed title style={boxrule=0pt, arc=2pt, colback=mcAmber}
}

\newtcolorbox{specbox}[1][Especificação]{
  enhanced, breakable, colback=white, colframe=mcLine,
  boxrule=0.8pt, left=8pt, right=8pt, top=6pt, bottom=6pt, arc=2pt,
  title={\iChip\ #1}, fonttitle=\bfseries\color{mcTealDark},
  coltitle=mcTealDark, colbacktitle=mcLight,
  attach boxed title to top left={yshift=-2mm, xshift=4mm},
  boxed title style={boxrule=0.8pt,colframe=mcLine, arc=2pt, colback=mcLight}
}

\newtcolorbox{resultbox}[1][Resultado]{
  enhanced, breakable, colback=mcLight, colframe=mcMint,
  boxrule=0.8pt, left=8pt, right=8pt, top=6pt, bottom=6pt, arc=2pt,
  title={\iCheck\ #1}, fonttitle=\bfseries\color{white},
  coltitle=white, colbacktitle=mcMint,
  attach boxed title to top left={yshift=-2mm, xshift=4mm},
  boxed title style={boxrule=0pt, arc=2pt, colback=mcMint}
}

% ------------------------------------------------------------
% Código-fonte (estilo listing)
% ------------------------------------------------------------
\lstdefinestyle{arduino}{
  backgroundcolor=\color{mcLight},
  basicstyle=\ttfamily\small\color{mcGray},
  keywordstyle=\color{mcTealDark}\bfseries,
  commentstyle=\color{mcMint}\itshape,
  stringstyle=\color{mcCoral},
  numbers=left, numberstyle=\tiny\color{mcGray},
  frame=single, rulecolor=\color{mcLine},
  breaklines=true, tabsize=2, showstringspaces=false,
  literate=
    {á}{{\'a}}1 {à}{{\`a}}1 {ã}{{\~a}}1 {â}{{\^a}}1
    {é}{{\'e}}1 {ê}{{\^e}}1 {í}{{\'i}}1 {ó}{{\'o}}1
    {õ}{{\~o}}1 {ô}{{\^o}}1 {ú}{{\'u}}1 {ç}{{\c c}}1
    {Á}{{\'A}}1 {À}{{\`A}}1 {Ã}{{\~A}}1 {Â}{{\^A}}1
    {É}{{\'E}}1 {Ê}{{\^E}}1 {Í}{{\'I}}1 {Ó}{{\'O}}1
    {Õ}{{\~O}}1 {Ô}{{\^O}}1 {Ú}{{\'U}}1 {Ç}{{\c C}}1
}

% ------------------------------------------------------------
% Placeholder de imagem (para o time inserir fotos depois)
% ------------------------------------------------------------
\newcommand{\imgplaceholder}[2][8cm]{%
  \begin{center}
  \begin{tikzpicture}
    \node[draw=mcLine, line width=1pt, dashed, fill=mcLight,
          minimum width=#1, minimum height=3.4cm, rounded corners=4pt] (box) {};
    \node[color=mcGray, align=center, font=\small] at (box) {\iImage\ \\[4pt] #2 \\ \textit{(inserir imagem aqui)}};
  \end{tikzpicture}
  \end{center}
}

% ------------------------------------------------------------
% Slot de imagem "inteligente": já vem com o nome do arquivo
% esperado (dentro da pasta img/). Basta subir o arquivo com
% esse nome exato no Overleaf que a imagem aparece sozinha;
% enquanto o arquivo não existir, mostra uma caixa indicando
% qual nome usar.
%
% CORREÇÃO: placement fixado em [H] (via pacote float) para que
% cada imagem fique exatamente onde foi inserida no texto, sem
% "flutuar" para outras seções/páginas.
% ------------------------------------------------------------
\newcommand{\projimg}[3][0.85\textwidth]{%
  \begin{figure}[H]
  \centering
  \IfFileExists{img/#2}{%
    \includegraphics[width=#1]{img/#2}%
    \caption{#3}%
  }{%
    \begin{tikzpicture}
      \node[draw=mcLine, line width=1pt, dashed, fill=mcLight,
            minimum width=#1, minimum height=3.4cm, rounded corners=4pt] (box) {};
      \node[color=mcGray, align=center, font=\small] at (box)
        {\iImage\ \\[4pt] #3 \\ \textit{arquivo esperado:} \texttt{img/#2}};
    \end{tikzpicture}%
  }%
  \end{figure}%
  \FloatBarrier
}

\newcolumntype{Y}{>{\raggedright\arraybackslash}X}

% ============================================================
\begin{document}

% ---------------- CAPA (estilo folha de produto) ----------------
\begin{titlepage}
\newgeometry{margin=0cm}
\begin{tikzpicture}[remember picture, overlay]
  \path (current page.north west) -- (current page.north east)
        coordinate[pos=0] (nw) coordinate[pos=1] (ne);
  \fill[mcTeal] (nw) rectangle ($(ne)+(0,-9.5cm)$);
  \fill[mcTealDark] ($(nw)+(0,-9.25cm)$) rectangle ($(ne)+(0,-9.5cm)$);

  % ícone tipo "comprimido / cápsula" no topo
  \node[anchor=north, yshift=-1.6cm] at (current page.north) {
    \begin{tikzpicture}
      \draw[fill=white, draw=none] (-0.9,0) -- (0.9,0) arc (-90:90:0.55) -- (-0.9,1.1) arc (90:270:0.55) -- cycle;
      \draw[fill=mcMint, draw=none] (0,-0.55) rectangle (0.9,0.65) ;
      \clip (-0.9,0) -- (0.9,0) arc (-90:90:0.55) -- (-0.9,1.1) arc (90:270:0.55) -- cycle;
      \fill[mcMint] (0,-0.6) rectangle (1,1.2);
    \end{tikzpicture}
  };

  \node[anchor=north, yshift=-3.6cm, text=white, font=\bfseries\fontsize{34}{34}\selectfont] at (current page.north) {MedControl};
  \node[anchor=north, yshift=-4.6cm, text=mcLight, font=\Large] at (current page.north) {Farmácia Inteligente};
  \node[anchor=north, yshift=-5.35cm, text=white!85, font=\normalsize] at (current page.north) {\iChip\ Sistema embarcado para administração automatizada de medicamentos};

  \node[anchor=north, yshift=-7.1cm, draw=white, line width=1pt, rounded corners=2pt, inner sep=6pt, text=white, font=\small\bfseries] at (current page.north) {DOCUMENTAÇÃO TÉCNICA \ \textbullet\ VERSÃO 1.0};
\end{tikzpicture}

\vspace*{10.6cm}

\begin{center}
{\LARGE\bfseries\color{mcTealDark} Relatório de Projeto}\\[0.3cm]
{\large\color{mcGray} Sistemas Embarcados aplicados à Saúde}
\end{center}

\vspace{1.0cm}

\begin{center}
\renewcommand{\arraystretch}{1.5}
\begin{tabular}{@{} l l @{}}
\textbf{\color{mcTealDark}Produto:} & MedControl -- Protótipo de Farmácia Inteligente \\
\textbf{\color{mcTealDark}Plataforma:} & Arduino Uno R3 \\
\textbf{\color{mcTealDark}Prototipagem:} & Tinkercad (simulação) + Maquete física \\
\textbf{\color{mcTealDark}Disciplina:} & Técnicas de Programação \\
\textbf{\color{mcTealDark}Semestre:} & 2026.1 \\
\textbf{\color{mcTealDark}Professor(a):} & Prof. Fabrício da Costa Dias \\
\textbf{\color{mcTealDark}Data de entrega:} & 21/07/2026 \\
\end{tabular}
\end{center}

\vspace{0.cm}

\begin{center}
\renewcommand{\arraystretch}{1.35}
\begin{tabular}{@{} l l @{}}
\multirow{7}{*}{\textbf{\color{mcTealDark}Equipe:}}
 & Conceição Tayná Medeiros da Silva \\
 & Gislane Araújo Silva \\
 & Luan Lopes dos Santos \\
 & Rafael Esdras Nascimento Soares \\
 & Viviane Freire Gomes \\
 & Wendell Gomes dos Reis \\
 & Yago Vinícius dos Santos Holanda \\
\end{tabular}
\end{center}

\vspace{2.0 cm}

\begin{center}
\color{mcGray}\small \iGit\ Repositório do código-fonte disponível na Seção 7
\end{center}

\restoregeometry
\end{titlepage}

% ---------------- SUMÁRIO ----------------
\tableofcontents
\newpage

% ============================================================
\section{Introdução}
% ============================================================

O \textbf{MedControl} é um protótipo de farmácia inteligente desenvolvido com Arduino, projetado para auxiliar usuários na administração correta de medicamentos. O sistema aciona a abertura do compartimento de medicação por meio de um servomotor e emite alertas sonoros e visuais para lembrar o usuário do momento da dose, aguardando a confirmação manual por botão antes de fechar o compartimento novamente.

Este documento apresenta a especificação técnica, a arquitetura e os resultados do desenvolvimento do produto, seguindo um formato de documentação técnica para facilitar consulta e manutenção futura.

\subsection{Objetivo Geral}

\begin{infobox}[Objetivo do Produto]
Desenvolver um protótipo de farmácia inteligente utilizando Arduino, capaz de auxiliar usuários na administração correta de medicamentos por meio da automação da abertura do compartimento e da emissão de alertas sonoros e visuais.
\end{infobox}

\subsection{Objetivos Específicos}

\begin{itemize}[leftmargin=1.4em, itemsep=2pt]
  \item Aplicar conceitos de sistemas embarcados no desenvolvimento de um produto funcional.
  \item Desenvolver um sistema automatizado para alerta e liberação de medicamentos.
  \item Integrar sensores e atuadores em uma única solução coesa.
  \item Construir uma maquete física utilizando materiais recicláveis.
  \item Simular o funcionamento completo no Tinkercad antes da implementação física.
  \item Demonstrar uma aplicação prática da automação na área da saúde.
\end{itemize}

\subsection{Escopo do Produto}

\begin{specbox}[Visão Geral]
\begin{tabularx}{\linewidth}{@{}l Y@{}}
\textbf{Nome do produto} & MedControl \\
\textbf{Categoria} & Dispositivo embarcado de assistência à saúde \\
\textbf{Público-alvo} & Pacientes com tratamentos de múltiplos horários, idosos, cuidadores \\
\textbf{Plataforma de hardware} & Arduino Uno R3 \\
\textbf{Ambiente de prototipagem} & Tinkercad (simulação) \\
\textbf{Status} & Protótipo funcional / maquete \\
\end{tabularx}
\end{specbox}

\newpage
% ============================================================
\section{Materiais e Equipamentos Utilizados}
% ============================================================

Esta seção documenta todos os componentes de hardware e materiais estruturais utilizados na construção do MedControl, de acordo com a lista de materiais (BOM) da simulação final no Tinkercad.

\subsection{Componentes Eletrônicos (Hardware)}

\begin{table}[H]
\centering
\begin{tabularx}{\linewidth}{@{}l Y c@{}}
\toprule
\textbf{Componente} & \textbf{Função no sistema} & \textbf{Qtd.} \\
\midrule
Arduino Uno R3 & Unidade central de processamento & 1 \\
Display LCD 16x2 & Exibição de informações ao usuário & 1 \\
Potenciômetro 5~k$\Omega$ & Ajuste de contraste do LCD & 1 \\
Servomotor (posicional) & Abertura/fechamento do compartimento & 1 \\
Piezo (buzzer) & Alerta sonoro & 1 \\
LED vermelho & Sinalização visual do alerta & 1 \\
Resistor 220~$\Omega$ & Proteção do LED / circuito & 2 \\
Push Button & Confirmação de dose tomada & 1 \\
Protoboard & Montagem do circuito & 1 \\
Jumpers & Conexões elétricas & — \\
\bottomrule
\end{tabularx}
\caption{Lista de componentes eletrônicos do MedControl (conforme BOM da simulação no Tinkercad).}
\end{table}

\begin{infobox}[Nota sobre o hardware]
Na simulação, o alerta sonoro é reproduzido por um \textbf{piezo}. Na montagem física do protótipo, esse mesmo componente é substituído por um \textbf{buzzer}, mantendo a mesma função e a mesma lógica de acionamento no firmware. Todos os demais componentes e conexões são idênticos entre a simulação e a maquete física.
\end{infobox}

\subsection{Materiais da Estrutura Física}

\begin{itemize}[leftmargin=1.4em, itemsep=2pt]
  \item Papelão reciclado
  \item Cola quente
  \item Fita adesiva
  \item Palitos de churrasco
  \item Papel contact
\end{itemize}

\begin{infobox}[Sustentabilidade]
A estrutura da maquete prioriza materiais recicláveis e de baixo custo, reforçando a viabilidade do protótipo e a consciência ambiental no processo de prototipagem.
\end{infobox}

\FloatBarrier
\newpage
% ============================================================
\section{Projeto no Tinkercad}
% ============================================================

\subsection{Simulação no Tinkercad}

O circuito completo foi projetado e validado na plataforma Tinkercad antes da montagem física, permitindo testar a lógica de controle sem riscos de dano aos componentes. A simulação final é composta por um Arduino Uno R3, um display LCD 16x2 (com potenciômetro para ajuste de contraste), um servomotor, um LED indicador, um piezo (buzzer) e um botão de confirmação, todos montados em protoboard.

\begin{figure}[H]
\centering
\IfFileExists{img/montagem_tinkercad.png}{%
  \includegraphics[width=0.92\textwidth]{img/montagem_tinkercad.png}%
}{%
  \begin{tikzpicture}
    \node[draw=mcLine, line width=1pt, dashed, fill=mcLight,
          minimum width=0.92\textwidth, minimum height=4cm, rounded corners=4pt] (box) {};
    \node[color=mcGray, align=center, font=\small] at (box)
      {\iImage\ \\[4pt] Montagem eletrônica do MedControl no Tinkercad \\ \textit{arquivo esperado:} \texttt{img/montagem\_tinkercad.png}};
  \end{tikzpicture}%
}
\caption{Montagem eletrônica do MedControl simulada no Tinkercad.}
\end{figure}
\FloatBarrier

\begin{figure}[H]
\centering
\IfFileExists{img/diagrama_esquematico.png}{%
  \includegraphics[width=0.92\textwidth]{img/diagrama_esquematico.png}%
}{%
  \begin{tikzpicture}
    \node[draw=mcLine, line width=1pt, dashed, fill=mcLight,
          minimum width=0.92\textwidth, minimum height=4cm, rounded corners=4pt] (box) {};
    \node[color=mcGray, align=center, font=\small] at (box)
      {\iImage\ \\[4pt] Diagrama esquemático do circuito completo \\ \textit{arquivo esperado:} \texttt{img/diagrama\_esquematico.png}};
  \end{tikzpicture}%
}
\caption{Diagrama esquemático do circuito completo do MedControl.}
\end{figure}
\FloatBarrier

\subsection{Mapeamento de Pinos}

\begin{table}[H]
\centering
\begin{tabularx}{0.8\linewidth}{@{}Y c@{}}
\toprule
\textbf{Componente} & \textbf{Pino} \\
\midrule
Servomotor (sinal) & D9 \\
Piezo / Buzzer & D10 \\
LED indicador & D8 \\
Botão de confirmação & D7 \\
LCD -- RS & D12 \\
LCD -- Enable (E) & D11 \\
LCD -- D4 & D5 \\
LCD -- D5 & D4 \\
LCD -- D6 & D3 \\
LCD -- D7 & D2 \\
\bottomrule
\end{tabularx}
\caption{Mapeamento de pinos do Arduino Uno.}
\end{table}

\subsection{Fluxo de Funcionamento}

O diagrama a seguir resume a lógica de operação do sistema, desde o disparo do alerta até a confirmação manual da dose pelo usuário.

\begin{figure}[H]
\centering
\resizebox{!}{0.62\textheight}{
\begin{tikzpicture}[
  node distance=0.4cm and 1cm,
  every node/.style={font=\small},
  proc/.style={rectangle, rounded corners=3pt, draw=mcTealDark, fill=mcLight, text=mcGray, minimum width=4.8cm, minimum height=0.9cm, align=center, drop shadow},
  dec/.style={diamond, draw=mcAmber, fill=white, text=mcGray, aspect=2.4, minimum width=4.8cm, align=center, drop shadow},
  startstop/.style={rectangle, rounded corners=12pt, draw=mcTeal, fill=mcTeal, text=white, minimum width=3.4cm, minimum height=0.8cm, align=center},
  arr/.style={-{Latex[length=2.2mm]}, draw=mcTealDark, thick}
]
  \node[startstop] (start) {Início};
  \node[proc, below=of start] (msginicial) {Exibir "MedControl" no LCD};
  \node[proc, below=of msginicial] (open) {Abrir compartimento (servo)};
  \node[proc, below=of open] (alert) {Acender LED + Tocar alerta sonoro};
  \node[proc, below=of alert] (lcd) {Mostrar "Aperte após tomar" no LCD};
  \node[dec, below=of lcd] (confirm) {Botão pressionado?};
  \node[proc, below=of confirm] (ok) {Mostrar "Confirmado!" e apagar alertas};
  \node[proc, below=of ok] (close) {Fechar compartimento (servo)};
  \node[proc, below=of close] (fim2) {Mostrar "Finalizado!"};
  \node[startstop, below=of fim2] (end) {Fim do ciclo};

  \draw[arr] (start) -- (msginicial);
  \draw[arr] (msginicial) -- (open);
  \draw[arr] (open) -- (alert);
  \draw[arr] (alert) -- (lcd);
  \draw[arr] (lcd) -- (confirm);
  \draw[arr] (confirm.east) .. controls +(2.4,0) and +(2.4,0) .. (alert.east) node[midway, right, font=\scriptsize\color{mcGray}] {não, aguarda};
  \draw[arr] (confirm) -- (ok) node[midway, right, xshift=2pt, font=\scriptsize\color{mcGray}] {sim};
  \draw[arr] (ok) -- (close);
  \draw[arr] (close) -- (fim2);
  \draw[arr] (fim2) -- (end);
\end{tikzpicture}
}
\caption{Fluxograma de funcionamento do MedControl.}
\end{figure}
\FloatBarrier

\newpage
% ============================================================
\section{Ilustrações/Imagens do Protótipo Físico}
% ============================================================

Registro visual da evolução do desenvolvimento, da concepção até o protótipo final.

\projimg{esboco-inicial.png}{Esboço/desenho inicial do projeto.}
\projimg{montagem-fisica.png}{Montagem física da maquete em construção.}
\projimg{prototipo-final.png}{Protótipo final montado.}

\FloatBarrier
\newpage
% ============================================================
\section{Dificuldades Encontradas}
% ============================================================

Durante o desenvolvimento do MedControl, a equipe enfrentou desafios técnicos comuns a projetos de sistemas embarcados:

\begin{warnbox}[Principais dificuldades]
\begin{itemize}[leftmargin=1.2em, itemsep=3pt]
  \item Sincronização entre o acionamento do servomotor e a abertura da tampa do compartimento.
  \item Ajuste do potenciômetro para garantir boa legibilidade do display LCD.
  \item Organização e gerenciamento dos fios dentro da estrutura da maquete.
  \item Dimensionamento adequado da estrutura física para acomodar os componentes.
  \item Adaptação da lógica do projeto ao ambiente de simulação do Tinkercad.
  \item Substituição do piezo (usado na simulação) pelo buzzer real na montagem física, mantendo a mesma lógica de firmware.
\end{itemize}
\end{warnbox}

\FloatBarrier
\newpage
% ============================================================

\section{Resultados Alcançados}
% ============================================================

\begin{resultbox}[Funcionamento do Protótipo]
O sistema foi capaz de acionar automaticamente o servomotor para abrir o compartimento, emitir alerta sonoro pelo piezo/buzzer, acionar sinalização luminosa pelo LED e apresentar informações no display LCD. O compartimento permanece aberto até que o usuário pressione o botão de confirmação, momento em que o LCD exibe a mensagem de confirmação e o servo retorna à posição fechada. Os testes realizados no Tinkercad demonstraram o funcionamento correto da lógica desenvolvida, validando o projeto antes da construção física.
\end{resultbox}

\subsection{Limitações Atuais}
\begin{itemize}[leftmargin=1.4em, itemsep=2pt]
  \item O sistema não possui controle automático de horário (RTC); o ciclo de alerta é disparado a cada repetição do loop principal.
  \item Ausência de persistência de dados (histórico de doses) em memória não-volátil.
  \item Estrutura física limitada a um único compartimento.
  \item Configuração remota (Wi-Fi) ainda não implementada.
\end{itemize}

\subsection{Melhorias Futuras}
\begin{itemize}[leftmargin=1.4em, itemsep=2pt]
  \item Adicionar módulo RTC para programação de horários reais de medicação.
  \item Implementar configuração via Wi-Fi (ex.: ESP8266) para ajuste remoto de horários.
  \item Adicionar aplicativo/mobile companion para acompanhamento remoto de doses.
  \item Implementar registro de histórico com cartão SD ou EEPROM.
  \item Expandir o número de compartimentos e adicionar sensores de presença do comprimido.
\end{itemize}

\FloatBarrier
\newpage
% ============================================================
\section{Código-Fonte do Projeto}
% ============================================================

O código-fonte completo do firmware desenvolvido para o Arduino encontra-se disponível no repositório abaixo.

\begin{specbox}[Repositório]
\iGit\ \ \url{https://github.com/rafaesdra/medcontrol} 
\end{specbox}

\subsection{Firmware Completo}

\begin{lstlisting}[style=arduino, language=C++, caption={Firmware do MedControl -- versão validada na simulação Tinkercad.}]
#include <LiquidCrystal.h>
#include <Servo.h>

// LCD: RS, E, D4, D5, D6, D7
LiquidCrystal lcd(12, 11, 5, 4, 3, 2);
Servo servo;

const int servoPin = 9;
const int buzzerPin = 10;
const int ledPin = 8;
const int botaoPin = 7;

// Notas
#define DO 262
#define RE 294
#define MI 330
#define FA 349
#define SOL 392
#define LA 440

void tocarAlerta() {
  int melodia[] = {
    DO, DO, SOL, SOL, LA, LA, SOL,
    FA, FA, MI, MI, RE, RE, DO
  };
  int duracao[] = {
    300, 300, 300, 300, 300, 300, 600,
    300, 300, 300, 300, 300, 300, 600
  };
  for (int i = 0; i < 14; i++) {
    tone(buzzerPin, melodia[i], duracao[i]);
    delay(duracao[i] + 50);
  }
  noTone(buzzerPin);
}

void setup() {
  servo.attach(servoPin);
  servo.write(0);
  pinMode(ledPin, OUTPUT);
  pinMode(buzzerPin, OUTPUT);
  pinMode(botaoPin, INPUT_PULLUP);
  digitalWrite(ledPin, LOW);

  lcd.begin(16, 2);
  lcd.setCursor(0, 0);
  lcd.print("MedControl");
  lcd.setCursor(0, 1);
  lcd.print("Iniciando...");
  delay(2000);
}

void loop() {
  // Abre compartimento
  lcd.clear();
  lcd.setCursor(0, 0);
  lcd.print("Hora remedio!");
  servo.write(90);
  digitalWrite(ledPin, HIGH);
  tocarAlerta();
  lcd.setCursor(0, 1);
  lcd.print("Aperte apos tomar");

  // Aguarda botao
  while (digitalRead(botaoPin) == HIGH) {
    // mantem o alerta piscando a cada alguns segundos
    delay(100);
  }

  // Confirmacao
  lcd.clear();
  lcd.setCursor(0, 0);
  lcd.print("Confirmado!");
  digitalWrite(ledPin, LOW);
  noTone(buzzerPin);
  delay(1000);

  // Fecha
  lcd.clear();
  lcd.setCursor(0, 0);
  lcd.print("Fechando...");
  servo.write(0);
  delay(1500);

  lcd.clear();
  lcd.setCursor(0, 0);
  lcd.print("Finalizado!");
  delay(5000);
}
\end{lstlisting}

\begin{infobox}[Nota]
O código acima corresponde à versão executada e validada na simulação do Tinkercad. Na montagem física do protótipo, o firmware é exatamente o mesmo; a única diferença de hardware é que o \textbf{piezo} usado na simulação é substituído por um \textbf{buzzer} real, sem qualquer alteração de código, pinagem ou lógica de funcionamento.
\end{infobox}

\newpage
% ============================================================
\section{Referências Bibliográficas}
% ============================================================

\begin{itemize}[leftmargin=1.4em, itemsep=4pt]
  \item ARDUINO. \textit{Documentação Oficial}. Disponível em: \url{https://docs.arduino.cc}.
  \item TINKERCAD. \textit{Circuits}. Disponível em: \url{https://www.tinkercad.com/things/ftH33SuNnAO-medcontrol?sharecode=Piq3hywgwHqZxrc1wp1MAs9D1p_18CyBY3xNzeBqxWE}.
  \item Datasheet do Arduino Uno R3.
  \item Datasheet do Servomotor (posicional).
  \item Datasheet do Display LCD 16x2.
  \item Datasheet do Buzzer/Piezo.
\end{itemize}

\vspace{2cm}
\begin{center}
\color{mcGray}\small
--- Fim da documentação técnica do MedControl ---
\end{center}

\end{document}
